\chapter{Finiteness descent from the residue field (Stacks 04GG fragment)}
\label{ch:Proetale_Mathlib_RingTheory_Module_FiniteResidue}

This chapter extracts a Mathlib-PR-quality helper for the
``finiteness descends from the residue field'' fragment of
Stacks Tag 04GG. The setting is a Noetherian local ring \(A\) with
maximal ideal \(\mathfrak{m}\) and residue field \(k = A/\mathfrak{m}\),
together with a finitely-presented flat \(A\)-module \(M\) whose
residue \(M \otimes_A k\) is finite-dimensional as a \(k\)-vector
space. The headline theorem is that \(M\) itself is module-finite
over \(A\).

The étale-specific corollary is the shape directly consumed by
\texttt{localisation\_finite\_at\_maximal} (see
\cref{ch:Proetale_Mathlib_RingTheory_Etale_LocalisationFiniteAtMaximal}).

% archon:covers Proetale/Mathlib/RingTheory/Module/FiniteResidue.lean

\section{Nakayama in residue form}

\begin{lemma}
  \label{lem:nakayama-residue-spanning}
  \lean{Module.eq_top_of_residue_eq_top_of_finite}
  [Nakayama from a residue spanning set]
  Let \(A\) be a local ring with maximal ideal \(\mathfrak{m}\), and
  let \(M\) be a finitely-generated \(A\)-module. If
  \(M = \mathfrak{m} \cdot M + N\) for some submodule
  \(N \subseteq M\), then \(N = M\).
\end{lemma}

\begin{proof}
  Apply Nakayama's lemma to the quotient \(M/N\). The hypothesis
  \(M = \mathfrak{m} \cdot M + N\) implies
  \(\mathfrak{m} \cdot (M/N) = M/N\): given any \(\bar x \in M/N\)
  with lift \(x \in M\), write \(x = m + n\) with \(m \in
  \mathfrak{m} \cdot M\) and \(n \in N\); then
  \(\bar x = \overline{m} \in \mathfrak{m} \cdot (M/N)\). Since
  \(M\) is finitely generated, so is the quotient \(M/N\). The
  maximal ideal \(\mathfrak{m}\) of a local ring is contained in
  its Jacobson radical, so Nakayama's lemma yields \(M/N = 0\),
  i.e.\ \(N = M\).

  \emph{Implementation note.} Mathlib may already provide this
  exact carrier under a name such as
  \texttt{Submodule.eq\_top\_of\_le\_smul\_of\_le\_jacobson\_of\_finite}
  or a Nakayama-as-iff variant; in that case the present lemma is a
  thin wrapper specialising the Jacobson hypothesis to the maximal
  ideal of a local ring. The wrapper is convenient because the
  downstream caller (\cref{thm:finite-of-finite-residue-flat-local})
  reasons about \(\mathfrak{m}\) directly.
\end{proof}

\section{Lifting a residue basis to a partial spanning set}

\begin{lemma}
  \label{lem:lift-residue-basis-to-spanning}
  \lean{Module.span_residue_basis_add_maximal_smul_eq_top}
  [Lift of residue basis to spanning modulo \(\mathfrak{m} M\)]
  Let \(A\) be a local ring with maximal ideal \(\mathfrak{m}\)
  and residue field \(k = A/\mathfrak{m}\), and let \(M\) be an
  \(A\)-module whose residue \(M \otimes_A k\) is finite-dimensional
  as a \(k\)-vector space, with basis
  \(\{\bar b_i\}_{i \in I}\). Choose any lift
  \(\{b_i\}_{i \in I} \subseteq M\) of the residue basis, in the
  sense that the image of \(b_i\) under the projection
  \(M \twoheadrightarrow M \otimes_A k\) (equivalently, under
  \(M \twoheadrightarrow M / \mathfrak{m} \cdot M\)) is \(\bar b_i\).
  Then
  \[
    \mathrm{span}_A \{b_i\}_{i \in I} \;+\; \mathfrak{m} \cdot M
    \;=\; M.
  \]
\end{lemma}

\begin{proof}
  Fix \(x \in M\) and consider its residue
  \(\bar x \in M \otimes_A k\). Under the canonical
  \(A\)-linear isomorphism
  \[
    M \otimes_A k \;\simeq\; M / \mathfrak{m} \cdot M,
  \]
  (the right-exactness of base change to \(A/\mathfrak{m}\), packaged
  by Mathlib as
  \texttt{Algebra.TensorProduct.quotIdealMapEquivQuotTensor} in the
  algebra case, or directly as the residue identification for
  modules), expand \(\bar x\) in the residue basis:
  \[
    \bar x \;=\; \sum_{i \in I} \alpha_i \bar b_i, \qquad \alpha_i \in k.
  \]
  Choose lifts \(a_i \in A\) of \(\alpha_i \in k = A/\mathfrak{m}\),
  for each \(i \in I\); since \(I\) is finite, only finitely many
  \(a_i\) are nonzero. Set
  \[
    y \;:=\; \sum_{i \in I} a_i \, b_i \;\in\; \mathrm{span}_A
      \{b_i\}_{i \in I} \;\subseteq\; M.
  \]
  The residue of \(y\) in \(M \otimes_A k\) is
  \(\sum_i \alpha_i \bar b_i = \bar x\) by construction of the
  lifts, hence \(x - y\) has residue \(0\). Under the identification
  \(M \otimes_A k \simeq M / \mathfrak{m} \cdot M\), this means
  \(x - y \in \mathfrak{m} \cdot M\). Therefore
  \[
    x \;=\; y \;+\; (x - y) \;\in\; \mathrm{span}_A \{b_i\}_{i \in I}
      \;+\; \mathfrak{m} \cdot M.
  \]
  Since \(x \in M\) was arbitrary, the displayed identity holds.
\end{proof}

\section{Headline: finiteness descent}

\begin{theorem}
  \label{thm:finite-of-finite-residue-flat-local}
  \lean{Module.finite_of_finite_residue_of_flat_of_local}
  [Stacks 04GG, finiteness descent fragment]
  Let \(A\) be a Noetherian local ring with maximal ideal
  \(\mathfrak{m}\) and residue field \(k\), and let \(M\) be a
  finitely-presented flat \(A\)-module. If \(M \otimes_A k\) is
  finite-dimensional as a \(k\)-vector space, then \(M\) is
  module-finite over \(A\).
\end{theorem}

\begin{proof}
  \uses{lem:lift-residue-basis-to-spanning, lem:nakayama-residue-spanning}
  \emph{Step 1: choose a residue basis and a lift.} Let
  \(d := \dim_k (M \otimes_A k)\), and pick a \(k\)-basis
  \(\{\bar b_i\}_{i \in I}\) of \(M \otimes_A k\), where \(I\) is a
  finite index set of cardinality \(d\). Choose lifts
  \(\{b_i\}_{i \in I} \subseteq M\) of the residue basis, i.e.\
  preimages under the projection
  \(M \twoheadrightarrow M / \mathfrak{m} \cdot M \simeq M \otimes_A k\).
  Let
  \[
    N \;:=\; \mathrm{span}_A \{b_i\}_{i \in I} \;\subseteq\; M.
  \]
  By construction, \(N\) is finitely generated.

  \emph{Step 2: residue surjectivity.} By
  \cref{lem:lift-residue-basis-to-spanning},
  \[
    N \;+\; \mathfrak{m} \cdot M \;=\; M.
  \]

  \emph{Step 3: \(M / N\) is finitely generated.} The module \(M\)
  is finitely-presented (by hypothesis), hence in particular
  finitely-generated. The quotient \(M/N\) is therefore the
  quotient of a finitely-generated \(A\)-module, hence itself
  finitely-generated. (Equivalently: a generating set of \(M\)
  descends to a generating set of \(M/N\).)

  \emph{Step 4: apply Nakayama.} Step~2 implies the projection
  \(M \twoheadrightarrow M/N\) sends \(M\) onto
  \(\mathfrak{m} \cdot (M/N)\), i.e.\
  \(\mathfrak{m} \cdot (M/N) = M/N\). Equivalently, treating
  \(N \subseteq M\) as the submodule of
  \cref{lem:nakayama-residue-spanning}, the hypothesis
  \(M = \mathfrak{m} \cdot M + N\) holds by Step~2.
  \cref{lem:nakayama-residue-spanning} applied to \(M\) and the
  submodule \(N\) (using Step~3 to know \(M\) is finitely generated)
  yields \(N = M\).

  \emph{Conclusion.} \(M = N = \mathrm{span}_A \{b_i\}_{i \in I}\),
  so the finite set \(\{b_i\}_{i \in I}\) generates \(M\) as an
  \(A\)-module. This is exactly \(\mathrm{Module.Finite}\,A\,M\).

  \emph{Remark on the flatness hypothesis.} The argument above
  uses only that \(M\) is finitely-presented (to obtain
  finite-generation of \(M\) and hence of \(M/N\) at Step~3), the
  residue identification \(M \otimes_A k \simeq M / \mathfrak{m} M\)
  (right-exactness of base change), and Nakayama. Flatness of \(M\)
  is the natural hypothesis in Stacks~04GG and is retained in the
  headline statement for direct alignment with downstream
  applications (where \(M\) is the underlying module of an étale
  \(A\)-algebra, hence automatically flat); the proof itself does
  not exploit flatness. A future refactor that drops the flatness
  hypothesis would not change the proof structure.
\end{proof}

\section{Étale specialisation}

\begin{corollary}
  \label{cor:finite-of-finite-residue-etale-local}
  \lean{Algebra.Etale.finite_of_finite_residue_of_local}
  [Stacks 04GG, étale fibre form]
  Let \(A\) be a Noetherian local ring with maximal ideal
  \(\mathfrak{m}\) and residue field \(k\), and let \(B\) be an
  étale \(A\)-algebra such that \(B \otimes_A k\) is finite-
  dimensional as a \(k\)-vector space. Then \(B\) is module-finite
  over \(A\).
\end{corollary}

\begin{proof}
  \uses{thm:finite-of-finite-residue-flat-local}
  An étale \(A\)-algebra is flat
  (\texttt{Algebra.Etale.flat}) and finitely-presented
  (\texttt{Algebra.Etale.finitePresentation}) over \(A\) by
  definition / standard Mathlib carriers. Applying
  \cref{thm:finite-of-finite-residue-flat-local} to the underlying
  \(A\)-module \(M := B\) — whose residue is the \(k\)-vector
  space \(B \otimes_A k\), finite-dimensional by hypothesis —
  yields \(\mathrm{Module.Finite}\,A\,B\).
\end{proof}

\section{References}

\begin{itemize}
\item Stacks Project, Tag~\href{https://stacks.math.columbia.edu/tag/04GG}{04GG} —
  the at-the-fibre / finiteness descent statement: a finitely
  presented flat module over a Noetherian local ring whose special
  fibre is finite-dimensional is module-finite. The primary source
  of \cref{thm:finite-of-finite-residue-flat-local} and
  \cref{cor:finite-of-finite-residue-etale-local}.
\item Stacks Project, Tag~\href{https://stacks.math.columbia.edu/tag/00DV}{00DV} —
  Nakayama's lemma in module form. The source of
  \cref{lem:nakayama-residue-spanning} (in the local-ring
  specialisation).
\item \Cref{ch:Proetale_Mathlib_RingTheory_Etale_LocalisationFiniteAtMaximal}
  — the downstream consumer: the étale-localised
  \(B' := B[1/e] \simeq B_{\mathfrak n}\) feeds
  \cref{cor:finite-of-finite-residue-etale-local} (with \(B'\) in
  place of \(B\)) to obtain the
  \(\mathrm{Module.Finite}\,A\,B_{\mathfrak n}\) producer required
  by Step~2 of \cref{thm:localisation-finite-at-maximal}.
\end{itemize}
